\documentclass{article}
\usepackage{../style/latex/dhruv_theory_style}

\newcommand{\vcdim}{\operatorname{VCdim}}
\newcommand{\growth}{\Gamma}
\newcommand{\classA}{\mathcal{H}_{2\mathrm{int}}}
\newcommand{\classB}{\mathcal{H}_{\mathrm{quad}}}
\newcommand{\zeroone}{\{0,1\}}

% -------------------------------------------------------
% Document metadata
% -------------------------------------------------------
\title{\textbf{DSC 190/291 Assignment 2}}
\author{Dhruv Patel}
\date{\today}

\begin{document}
\maketitle
\tableofcontents
\newpage

\section{Part A: Restriction Patterns for Two Intervals}

\subsection{Problem Overview}

In Part A we study the hypothesis class consisting of unions of at most two
intervals on the real line. Let
\[
\classA
=
\left\{
\mathbf{1}_{I_1 \cup I_2}
:
I_1,I_2 \text{ intervals in } \R
\right\},
\]
where either interval is allowed to be empty. When we restrict a hypothesis in
\(\classA\) to an ordered sample
\[
x_1 < x_2 < \cdots < x_n,
\]
we obtain a binary string
\[
\bigl(h(x_1),\dots,h(x_n)\bigr)\in \zeroone^n.
\]
The main goal of Part A is to understand exactly which binary strings can arise,
count them, determine the VC dimension, and compare the exact counts to general
bounds from learning theory.

\subsection{Problem A1: Restriction Patterns}

\textbf{Question.}
\newline
Characterize exactly which binary labelings on an ordered sample
\[
x_1 < x_2 < \cdots < x_n
\]
are realizable by a hypothesis in \(\classA\).

The key statement to prove is:
\[
\text{a labeling is realizable}
\quad\Longleftrightarrow\quad
\text{it has at most two runs of }1\text{s}.
\]

\begin{solutionbox}

\subsubsection{Restatement}

We are given \(n\) ordered points on the real line and a binary labeling
\[
y=(y_1,\dots,y_n)\in\zeroone^n.
\]
We must determine exactly when there exists a union of at most two intervals
whose indicator function agrees with this labeling on the sample.

\subsubsection{Explanation}

An interval intersects an ordered sample in a contiguous block of indices. Since
our class allows at most two intervals, the positive labels should form at most
two contiguous blocks. In binary-string language, contiguous blocks of positive
labels are exactly runs of \(1\)s. So the whole problem is to make precise the
statement that unions of two intervals correspond exactly to binary strings with
at most two runs of \(1\)s.

\subsubsection{Solution Plan}

\begin{enumerate}[leftmargin=2.5em]
    \item Prove the forward implication:
    every hypothesis in \(\classA\) induces at most two runs of \(1\)s.
    \item Prove the reverse implication:
    every binary string with at most two runs of \(1\)s can be realized by
    explicitly constructing suitable intervals.
    \item Conclude the equivalence.
\end{enumerate}

\subsubsection{Formal Solution / Proof}

\begin{theorem}
Let \(x_1 < \cdots < x_n\) be distinct real numbers. A labeling
\[
(y_1,\dots,y_n)\in\zeroone^n
\]
is realizable by \(\classA\) if and only if the binary string
\[
y_1y_2\cdots y_n
\]
contains at most two runs of \(1\)s.
\end{theorem}

\begin{proof}
\leavevmode\\
\textbf{(\(\Rightarrow\)) Suppose \(h\in\classA\).}

Then there exist intervals \(I_1\) and \(I_2\) such that
\[
h=\mathbf{1}_{I_1\cup I_2}.
\]
For a sorted sample, the set of indices \(i\) with \(x_i\in I_1\) is either
empty or a contiguous block. Indeed, once the sample enters an interval, it
cannot leave and later re-enter that same interval while the sample points are
ordered increasingly. The same statement is true for \(I_2\).

Therefore the set of indices with label \(1\) under \(h\) is the union of at
most two contiguous blocks. Equivalently, the binary string
\[
h(x_1)\cdots h(x_n)
\]
has at most two runs of \(1\)s.

\textbf{(\(\Leftarrow\)) Suppose the string \(y_1\cdots y_n\) has at most two
runs of \(1\)s.}

We construct a hypothesis in \(\classA\) that realizes it.

\textbf{Case 1: no runs of \(1\)s.}

Then every label is \(0\). This is realized by taking the union of two empty
intervals.

\textbf{Case 2: exactly one run of \(1\)s.}

Suppose the run occupies indices \(a,\dots,b\). Then define one interval \(I_1\)
that contains exactly the sample points \(x_a,\dots,x_b\). For example, if
\(1<a\) and \(b<n\), we may take
\[
I_1=
\left[\frac{x_{a-1}+x_a}{2},\,\frac{x_b+x_{b+1}}{2}\right].
\]
If the run touches the left or right edge of the sample, we modify the
corresponding endpoint in the obvious way. Let \(I_2\) be empty. Then
\[
\mathbf{1}_{I_1\cup I_2}(x_i)=y_i
\qquad\text{for all }i.
\]

\textbf{Case 3: exactly two runs of \(1\)s.}

Suppose the two runs occupy blocks \(a,\dots,b\) and \(c,\dots,d\), where
\[
1\le a\le b < c\le d\le n.
\]
Choose one interval \(I_1\) containing exactly the first block and another
interval \(I_2\) containing exactly the second block. Then
\[
\mathbf{1}_{I_1\cup I_2}(x_i)=y_i
\qquad\text{for all }i.
\]

Thus every binary string with at most two runs of \(1\)s is realizable by a
union of at most two intervals.
\end{proof}

\subsubsection{Conclusion}

The realizable restriction patterns of \(\classA\) are exactly the binary
strings with at most two runs of \(1\)s.

\subsubsection{Lean Formalization}

Lean: NO.

\subsubsection{Computational Check}

For \(n\le 6\), the script \texttt{hw2/code/enumeration\_check.py} checks that
the counting formula from Problem A2 matches the number of binary strings with
at most two runs of \(1\)s.

\end{solutionbox}

\subsection{Problem A3: VC Dimension and Halving Bound}

\textbf{Question.}
\newline
Determine the VC dimension of \(\classA\). Show that \(n=4\) is shatterable and
that \(n=5\) is not. Then state the exact Halving bound and compare it to the
Sauer--Shelah bound.

\begin{solutionbox}

\subsubsection{Restatement}

We want the largest sample size on which \(\classA\) can realize every binary
labeling. Then we use the exact growth function from Problem A2 to state the
Halving mistake bound.

\subsubsection{Explanation}

The run characterization from Problem A1 makes the VC-dimension argument very
clean. Every binary string on four bits has at most two runs of \(1\)s, but the
five-bit string \(10101\) has three runs of \(1\)s and therefore is not
realisable.

\subsubsection{Solution Plan}

\begin{enumerate}[leftmargin=2.5em]
    \item Show that every labeling of four ordered points is realizable.
    \item Exhibit a nonrealizable labeling on five ordered points.
    \item Conclude \(\vcdim(\classA)=4\).
    \item Substitute the exact growth function into \(\log_2 \growth_{\classA}(n)\).
\end{enumerate}

\subsubsection{Formal Solution / Proof}

\begin{claim}
\[
\vcdim(\classA)=4.
\]
\end{claim}

\begin{proof}
\leavevmode\\
Take any four ordered points
\[
x_1<x_2<x_3<x_4.
\]
Any binary string of length \(4\) has at most two runs of \(1\)s. Therefore by
Problem A1 every labeling on these four points is realizable by \(\classA\). So
\(\{x_1,x_2,x_3,x_4\}\) is shattered and
\[
\vcdim(\classA)\ge 4.
\]

Now take any five ordered points
\[
x_1<x_2<x_3<x_4<x_5.
\]
The labeling
\[
(1,0,1,0,1)
\]
has three runs of \(1\)s. By Problem A1 this labeling is not realizable by a
union of at most two intervals. Therefore no \(5\)-point set is shattered, so
\[
\vcdim(\classA)\le 4.
\]
Hence
\[
\vcdim(\classA)=4.
\]
\end{proof}

From Problem A2,
\[
\growth_{\classA}(n)=1+\binom{n+1}{2}+\binom{n+1}{4}.
\]
Therefore the exact Halving bound is
\[
M_{\mathrm{Halving}}
\le
\log_2\!\left(1+\binom{n+1}{2}+\binom{n+1}{4}\right).
\]

Since \(\vcdim(\classA)=4\), Sauer--Shelah yields
\[
\growth_{\classA}(n)\le \sum_{k=0}^4 \binom{n}{k},
\]
so the generic VC-based Halving bound is
\[
M_{\mathrm{Halving}}
\le
\log_2\!\left(\sum_{k=0}^4 \binom{n}{k}\right).
\]

\subsubsection{Conclusion}

The VC dimension is exactly
\[
\vcdim(\classA)=4,
\]
and the exact Halving bound is sharper than the generic Sauer--Shelah-based
bound because it uses the full structure of the class.

\subsubsection{Lean Formalization}

Lean: NO.

\subsubsection{Computational Check}

The exact growth values satisfy
\[
\growth_{\classA}(4)=16=2^4,
\qquad
\growth_{\classA}(5)=31<32=2^5,
\]
which agrees with the VC-dimension conclusion.

\end{solutionbox}

\subsection{Problem A4: Tightness of Sauer--Shelah}

\textbf{Question.}
\newline
Construct a hypothesis class with VC dimension exactly \(d\) and growth function
\[
\sum_{k=0}^{d}\binom{n}{k}.
\]
Explain why this shows the Sauer--Shelah bound can be tight.

\begin{solutionbox}

\subsubsection{Restatement}

We need an explicit class whose growth function exactly equals the
Sauer--Shelah upper bound.

\subsubsection{Explanation}

The natural example is the class of indicator functions of subsets of size at
most \(d\). On an \(n\)-point sample, the realizable labelings are exactly those
with at most \(d\) ones.

\subsubsection{Solution Plan}

\begin{enumerate}[leftmargin=2.5em]
    \item Define \(\mathcal{H}_{\le d}\).
    \item Show that every \(d\)-point set is shattered.
    \item Show that no \((d+1)\)-point set is shattered.
    \item Count the realizable labelings on an \(n\)-point sample.
\end{enumerate}

\subsubsection{Formal Solution / Proof}

Let \(X\) be an infinite set and define
\[
\mathcal{H}_{\le d}
=
\{\mathbf{1}_S : S\subseteq X,\ |S|\le d\}.
\]

\begin{claim}
\[
\vcdim(\mathcal{H}_{\le d})=d.
\]
\end{claim}

\begin{proof}
\leavevmode\\
Take any \(d\) distinct points \(x_1,\dots,x_d\in X\). For any labeling
\((y_1,\dots,y_d)\in\zeroone^d\), let
\[
S=\{x_i:y_i=1\}.
\]
Then \(|S|\le d\), so \(\mathbf{1}_S\in\mathcal{H}_{\le d}\), and it realizes
the labeling. Hence every \(d\)-point set is shattered, so
\[
\vcdim(\mathcal{H}_{\le d})\ge d.
\]

Now consider \(d+1\) distinct points \(x_1,\dots,x_{d+1}\). The all-ones
labeling on these points would require a set \(S\) containing all \(d+1\)
points, which is impossible because hypotheses in \(\mathcal{H}_{\le d}\) have
support size at most \(d\). Therefore no \((d+1)\)-point set is shattered, so
\[
\vcdim(\mathcal{H}_{\le d})\le d.
\]
Hence
\[
\vcdim(\mathcal{H}_{\le d})=d.
\]
\end{proof}

On an \(n\)-point sample, a labeling is realizable exactly when it has at most
\(d\) ones. For each \(k\), there are \(\binom{n}{k}\) labelings with exactly
\(k\) ones, so
\[
\growth_{\mathcal{H}_{\le d}}(n)=\sum_{k=0}^{d}\binom{n}{k}.
\]

\subsubsection{Conclusion}

This class has VC dimension exactly \(d\) and achieves the Sauer--Shelah upper
bound exactly. Therefore the inequality can be tight.

\subsubsection{Lean Formalization}

Lean: NO.

\subsubsection{Computational Check}

Not needed beyond the direct combinatorial count.

\end{solutionbox}

\subsection{Problem A5: AI Proof Audit}

\textbf{Question.}
\newline
An AI proof counts possible switch points as if a union of two intervals could
realize any binary string with up to four changes. Explain why this is invalid
and give the correct statement.

\begin{solutionbox}

\subsubsection{Restatement}

We must identify the flaw in a switch-point counting argument and replace it
with the correct run-based reasoning.

\subsubsection{Explanation}

The error is that switch points are not free parameters. The geometry of two
intervals forces the positive set to be the union of at most two contiguous
blocks, which is a stronger condition than merely bounding the number of label
changes.

\subsubsection{Solution Plan}

\begin{enumerate}[leftmargin=2.5em]
    \item Explain why arbitrary change-point counting overcounts.
    \item State the correct two-run characterization.
    \item Deduce the correct growth-function formula.
\end{enumerate}

\subsubsection{Formal Solution / Proof}

The AI proof is invalid because it ignores interval structure. A union of two
intervals does not produce an arbitrary collection of sign changes; it produces
at most two contiguous positive blocks. Therefore the correct structural
statement is:
\[
\text{a labeling is realizable by }\classA
\quad\Longleftrightarrow\quad
\text{it has at most two runs of }1\text{s}.
\]
Once this is proved, the exact growth function follows from Problem A2:
\[
\growth_{\classA}(n)=1+\binom{n+1}{2}+\binom{n+1}{4}.
\]

\subsubsection{Conclusion}

The AI argument overcounts because it counts unconstrained change patterns. The
correct argument is based on runs of \(1\)s.

\subsubsection{Lean Formalization}

Lean: NO.

\subsubsection{Computational Check}

The small-\(n\) enumeration supports the corrected formula.

\end{solutionbox}

\section{Part B: Quadratic Thresholds}

\subsection{Problem Overview}

In Part B we study
\[
\classB
=
\left\{
x\mapsto \mathbf{1}[a x^2+b x+c\ge 0]
:
(a,b,c)\in\R^3\setminus\{(0,0,0)\}
\right\}.
\]
The main tasks are to represent this class linearly, describe its restriction
patterns, count them, and compute its VC dimension.

\subsection{Problem B1: General Upper Bound for Linear Representations}

\textbf{Question.}
\newline
Prove that if a hypothesis class is represented as homogeneous linear threshold
functions in \(\R^D\), then its VC dimension is at most \(D\).

\begin{solutionbox}

\subsubsection{Restatement}

Suppose
\[
\mathcal{H}_{\phi}
=
\left\{
x\mapsto \mathbf{1}[\langle w,\phi(x)\rangle\ge 0]
:
w\in\R^D
\right\}.
\]
We must prove
\[
\vcdim(\mathcal{H}_{\phi})\le D.
\]

\subsubsection{Explanation}

The proof uses linear dependence. If \(D+1\) points were shattered, then their
feature vectors would realize every binary labeling by a homogeneous hyperplane.
But \(D+1\) vectors in \(\R^D\) are linearly dependent, and that dependence
produces a labeling that no such hyperplane can realize.

\subsubsection{Solution Plan}

\begin{enumerate}[leftmargin=2.5em]
    \item Assume \(D+1\) points are shattered.
    \item Write a nontrivial dependence relation among their feature vectors.
    \item Convert coefficient signs into a target labeling.
    \item Derive a contradiction by taking an inner product with the separating
    weight vector.
\end{enumerate}

\subsubsection{Formal Solution / Proof}

\begin{theorem}
Let \(X\) be any domain and \(\phi:X\to\R^D\). Define
\[
\mathcal{H}_{\phi}
=
\left\{
x\mapsto \mathbf{1}[\langle w,\phi(x)\rangle\ge 0]
:
w\in\R^D
\right\}.
\]
Then
\[
\vcdim(\mathcal{H}_{\phi})\le D.
\]
\end{theorem}

\begin{proof}
\leavevmode\\
Assume for contradiction that \(\mathcal{H}_{\phi}\) shatters
\[
x_1,\dots,x_{D+1}\in X.
\]
Let
\[
v_i=\phi(x_i)\in\R^D.
\]
Since there are \(D+1\) vectors in \(\R^D\), they are linearly dependent. So
there exist scalars \(\alpha_1,\dots,\alpha_{D+1}\), not all zero, such that
\[
\sum_{i=1}^{D+1}\alpha_i v_i=0.
\]

Remove any zero coefficients. On the remaining indices define the labeling
\[
y_i=
\begin{cases}
1, & \alpha_i>0,\\
0, & \alpha_i<0.
\end{cases}
\]
Because the sample is shattered, there exists \(w\in\R^D\) such that
\[
\alpha_i>0 \implies \langle w,v_i\rangle >0,
\qquad
\alpha_i<0 \implies \langle w,v_i\rangle <0.
\]
Now take inner products with \(w\):
\[
0
=
\left\langle
w,\sum_{i=1}^{D+1}\alpha_i v_i
\right\rangle
=
\sum_{i=1}^{D+1}\alpha_i\langle w,v_i\rangle.
\]
But every term in this final sum is strictly positive, because \(\alpha_i\) and
\(\langle w,v_i\rangle\) have the same sign. Therefore the sum is strictly
positive, which is impossible. This contradiction proves
\[
\vcdim(\mathcal{H}_{\phi})\le D.
\]
\end{proof}

\subsubsection{Conclusion}

Any class represented by homogeneous linear thresholds in \(\R^D\) has VC
dimension at most \(D\).

\subsubsection{Lean Formalization}

Lean: YES.

Target theorem:
\[
\mathcal{H}_{\phi}
=
\{x\mapsto \mathbf{1}[\langle w,\phi(x)\rangle\ge 0] : w\in\R^D\}
\Longrightarrow
\vcdim(\mathcal{H}_{\phi})\le D.
\]

Repository file:
\texttt{hw2/lean/vc\_dimension\_bound.lean}.

\subsubsection{Computational Check}

Not applicable.

\end{solutionbox}

\subsection{Problem B2: Quadratic Transformation}

\textbf{Question.}
\newline
Define
\[
\phi(x)=(x^2,x,1).
\]
Show that quadratic thresholds become linear thresholds under this map and
conclude that \(\vcdim(\classB)\le 3\).

\begin{solutionbox}

\subsubsection{Restatement}

We must show that
\[
a x^2+b x+c
=
\langle (a,b,c),(x^2,x,1)\rangle
\]
and then apply Problem B1 with \(D=3\).

\subsubsection{Explanation}

This is the standard feature-lifting idea: the original class is nonlinear in
the scalar input \(x\), but linear in the lifted coordinates \((x^2,x,1)\).

\subsubsection{Solution Plan}

\begin{enumerate}[leftmargin=2.5em]
    \item Write the quadratic polynomial as a dot product in \(\R^3\).
    \item Observe that the threshold class is now a linear-threshold class.
    \item Apply the general VC-dimension upper bound from Problem B1.
\end{enumerate}

\subsubsection{Formal Solution / Proof}

Define
\[
\phi:\R\to\R^3,
\qquad
\phi(x)=(x^2,x,1).
\]
For \(a,b,c\in\R\), let \(w=(a,b,c)\in\R^3\). Then
\[
\langle w,\phi(x)\rangle
=
\langle (a,b,c),(x^2,x,1)\rangle
=
a x^2+b x+c.
\]
Therefore
\[
\mathbf{1}[a x^2+b x+c\ge 0]
=
\mathbf{1}[\langle w,\phi(x)\rangle \ge 0].
\]
So \(\classB\) is a homogeneous linear-threshold class in \(\R^3\). By Problem
B1,
\[
\vcdim(\classB)\le 3.
\]

\subsubsection{Conclusion}

The quadratic-threshold class lifts to a linear-threshold class in three
dimensions, so
\[
\vcdim(\classB)\le 3.
\]

\subsubsection{Lean Formalization}

Lean: YES.

Target identity:
\[
a x^2+b x+c=\langle (a,b,c),(x^2,x,1)\rangle.
\]

Repository file:
\texttt{hw2/lean/quadratic\_embedding.lean}.

\subsubsection{Computational Check}

Not needed for the algebraic identity itself.

\end{solutionbox}

\subsection{Problem B3: Restriction Patterns}

\textbf{Question.}
\newline
Describe all binary labelings on an ordered sample
\[
x_1<x_2<\cdots<x_n
\]
that can be realized by thresholding a quadratic polynomial.

\begin{solutionbox}

\subsubsection{Restatement}

We must characterize all strings of the form
\[
\bigl(\mathbf{1}[q(x_1)\ge 0],\dots,\mathbf{1}[q(x_n)\ge 0]\bigr),
\]
where \(q(x)=a x^2+b x+c\).

\subsubsection{Explanation}

The sign of a quadratic can only change at real roots, and a quadratic has at
most two real roots. Therefore the nonnegative set is either all of \(\R\),
empty, one interval, or the complement of one interval. On a sorted sample these
geometric possibilities become binary-pattern constraints.

\subsubsection{Solution Plan}

\begin{enumerate}[leftmargin=2.5em]
    \item Analyze sign patterns according to the number of roots and the sign of
    the leading coefficient.
    \item Translate each geometric case to a binary pattern on the sample.
    \item Prove the converse by constructing quadratics for each valid pattern.
\end{enumerate}

\subsubsection{Formal Solution / Proof}

\begin{theorem}
On an ordered sample \(x_1<\cdots<x_n\), the realizable \(1\)-sets of
\(\classB\) are exactly:
\begin{enumerate}[leftmargin=2.5em]
    \item the empty set,
    \item one contiguous block,
    \item the union of a prefix and a suffix, separated by a nonempty middle
    block of zeros.
\end{enumerate}
\end{theorem}

\begin{proof}
\leavevmode\\
Let \(q(x)=a x^2+b x+c\).

If \(a=0\), then \(q\) is linear or constant, so the set \(\{x:q(x)\ge 0\}\) is
either empty, all of \(\R\), a left half-line, or a right half-line. On an
ordered sample this yields all zeros, all ones, a prefix of ones, or a suffix
of ones.

Now suppose \(a\neq 0\). A quadratic has at most two real roots.

If there are no real roots, the sign never changes, so the labeling is either
all zeros or all ones.

If there are two distinct real roots \(r_1<r_2\), then:
\begin{itemize}[leftmargin=2.5em]
    \item if \(a<0\), then \(q(x)\ge 0\) exactly on \([r_1,r_2]\), which gives
    one contiguous block of ones on the sample;
    \item if \(a>0\), then \(q(x)\ge 0\) exactly outside \((r_1,r_2)\), which
    gives a prefix together with a suffix.
\end{itemize}

Thus every realizable pattern has one of the three stated forms.

Conversely, the empty labeling is realized by \(q(x)=-1\), and the all-ones
labeling by \(q(x)=1\).

If we want one contiguous block from \(x_a\) through \(x_b\), choose roots
\[
r_1<x_a<x_b<r_2
\]
with \(r_1\) and \(r_2\) placed in the neighboring gaps, and define
\[
q(x)=-(x-r_1)(x-r_2).
\]
Then \(q(x)\ge 0\) exactly on the chosen block.

If we want a prefix together with a suffix, choose an interior zero block from
\(x_a\) through \(x_b\), place roots around it, and define
\[
q(x)=(x-r_1)(x-r_2).
\]
Then \(q(x)\ge 0\) exactly outside the middle interval, so the sample labeling
is the desired prefix--suffix pattern.
\end{proof}

\subsubsection{Conclusion}

Quadratic thresholds realize exactly three kinds of positive sets on an ordered
sample: empty, one block, or prefix-plus-suffix.

\subsubsection{Lean Formalization}

Lean: NO.

\subsubsection{Computational Check}

The small-\(n\) enumeration checks that counting these three families produces
the exact growth formula in Problem B4.

\end{solutionbox}

\subsection{Problem B4: Growth Function and VC Dimension}

\textbf{Question.}
\newline
Using the restriction-pattern description from Problem B3, compute the exact
growth function of \(\classB\), determine its VC dimension, and compare the
exact count with the Sauer--Shelah bound.

\begin{solutionbox}

\subsubsection{Restatement}

We must count exactly how many labelings arise from the three valid pattern
types in Problem B3 and then use the result to determine the VC dimension.

\subsubsection{Explanation}

The counting is clean once the structure is correct. The main subtlety is that a
prefix--suffix labeling is determined by its middle zero block, which must be a
nonempty contiguous interior block.

\subsubsection{Solution Plan}

\begin{enumerate}[leftmargin=2.5em]
    \item Count the empty labeling.
    \item Count one contiguous block of ones.
    \item Count prefix--suffix patterns.
    \item Add the counts and simplify.
    \item Check the formula at \(n=3\) and \(n=4\) to get the VC dimension.
\end{enumerate}

\subsubsection{Formal Solution / Proof}

By Problem B3, realizable labelings fall into three disjoint classes.

\textbf{Case 1: the empty labeling.}

There is exactly one such labeling.

\textbf{Case 2: one contiguous block of \(1\)s.}

Such a labeling is specified by choosing indices
\[
1\le a\le b\le n
\]
and labeling exactly \(x_a,\dots,x_b\) by \(1\). Therefore the number of such
labelings is
\[
\binom{n+1}{2}.
\]

\textbf{Case 3: a prefix together with a suffix.}

Such a labeling is determined by a nonempty contiguous middle block of zeros
lying strictly inside the sample. If the zero block runs from index \(a\) to
index \(b\), then
\[
2\le a\le b\le n-1.
\]
The number of such contiguous blocks in \(\{2,\dots,n-1\}\) is
\[
\binom{n-1}{2}.
\]

Therefore
\[
\growth_{\classB}(n)
=
1+\binom{n+1}{2}+\binom{n-1}{2}.
\]
Simplifying,
\begin{align*}
\growth_{\classB}(n)
&=
1+\frac{n(n+1)}{2}+\frac{(n-1)(n-2)}{2}\\
&=
1+\frac{n^2+n+n^2-3n+2}{2}\\
&=
1+\frac{2n^2-2n+2}{2}\\
&=
n^2-n+2.
\end{align*}

Now compute the VC dimension. At \(n=3\),
\[
\growth_{\classB}(3)=3^2-3+2=8=2^3,
\]
so three points are shattered. At \(n=4\),
\[
\growth_{\classB}(4)=4^2-4+2=14<16=2^4,
\]
so four points are not shattered. Hence
\[
\vcdim(\classB)=3.
\]

Since \(\vcdim(\classB)=3\), Sauer--Shelah gives
\[
\growth_{\classB}(n)\le \sum_{k=0}^{3}\binom{n}{k}.
\]
The exact count \(n^2-n+2\) is sharper because it uses the actual structure of
quadratic thresholds.

\subsubsection{Conclusion}

The exact growth function is
\[
\growth_{\classB}(n)=n^2-n+2,
\]
and the VC dimension is exactly
\[
\vcdim(\classB)=3.
\]

\subsubsection{Lean Formalization}

Lean: OPTIONAL.

\subsubsection{Computational Check}

The script verifies the values
\[
2,\ 4,\ 8,\ 14,\ 22,\ 32
\]
for \(n=1,\dots,6\).

\end{solutionbox}

\subsection{Problem B5: AI Proof Audit}

\textbf{Question.}
\newline
An AI proof counts arbitrary change points and concludes a growth function for
quadratic thresholds. Explain the flaw and replace it with the correct
root-based reasoning.

\begin{solutionbox}

\subsubsection{Restatement}

We must diagnose why counting change points alone is mathematically insufficient
for quadratic thresholds.

\subsubsection{Explanation}

The issue is that not every pattern with a small number of sign changes comes
from a quadratic. The shape of the polynomial matters: the positive set must be
an interval or the complement of an interval.

\subsubsection{Solution Plan}

\begin{enumerate}[leftmargin=2.5em]
    \item Explain why free change-point counting overcounts.
    \item State the actual structural constraint from Problem B3.
    \item Deduce the correct growth function from Problem B4.
\end{enumerate}

\subsubsection{Formal Solution / Proof}

The AI proof is invalid because it counts change locations as if they could be
chosen independently. That ignores the geometry of quadratics. A quadratic
polynomial has at most two real roots, and depending on the sign of the leading
coefficient, the region where \(q(x)\ge 0\) is either one interval or the
complement of one interval.

Therefore the correct restriction patterns are not ``all strings with at most
two changes.'' The correct statement from Problem B3 is that the realizable
\(1\)-sets are exactly:
\begin{itemize}[leftmargin=2.5em]
    \item the empty set,
    \item one contiguous block,
    \item a prefix together with a suffix.
\end{itemize}
Using this classification, the correct growth function is
\[
\growth_{\classB}(n)=n^2-n+2.
\]

\subsubsection{Conclusion}

The flawed AI proof ignores the shape constraints imposed by quadratics. The
correct proof must be root-based rather than arbitrary-change-based.

\subsubsection{Lean Formalization}

Lean: NO.

\subsubsection{Computational Check}

The enumeration script supports the corrected formula on small values of \(n\).

\end{solutionbox}

\section{Part C: AI Usage Report}

\textbf{Question.}
\newline
Write a short report describing where AI was used, an accepted suggestion, a
rejected or modified suggestion, how correctness was verified, and how the
workflow improved.

\begin{solutionbox}

\subsubsection{Restatement}

This part asks for a reflective description of how AI contributed to the work
without replacing mathematical verification.

\subsubsection{Explanation}

The goal is to show judgment: where AI was helpful, where it needed correction,
and what independent checks were used to make sure the final results were
correct.

\subsubsection{Solution Plan}

\begin{enumerate}[leftmargin=2.5em]
    \item Describe where AI was used in the workflow.
    \item Give one accepted suggestion.
    \item Give one rejected or revised suggestion.
    \item Explain the verification process.
    \item Summarize workflow improvements.
\end{enumerate}

\subsubsection{Formal Solution / Proof}

For this assignment, I used AI mainly for organization, proof drafting, and
small computational support. In the proof-writing stage, AI helped turn rough
ideas into a cleaner structure with definitions, claims, and proofs. In the
verification stage, AI helped suggest a small brute-force enumeration script for
checking the growth-function formulas on small values of \(n\). I also used AI
to identify which parts of the assignment were the most natural targets for Lean
formalization.

One suggestion I accepted was the idea of using a small-\(n\) computational
check for the exact growth formulas. I accepted this because the formulas in
Parts A and B are exact, so enumeration on \(n\le 6\) is a meaningful sanity
check. It does not replace proof, but it helps catch overcounting and missed
cases.

One suggestion I rejected or heavily modified was an earlier proof sketch that
counted label changes too loosely. That approach risked overcounting because it
did not properly encode interval structure in Part A or root structure in Part
B. I replaced it with arguments based on runs of \(1\)s for the two-interval
class and on the sign pattern of quadratics for the quadratic-threshold class.

I verified correctness in several ways. First, every counting formula was
derived from a structural characterization proved immediately before it. Second,
I checked the formulas computationally on small samples:
\[
\growth_{\classA}(n)=1+\binom{n+1}{2}+\binom{n+1}{4},
\qquad
\growth_{\classB}(n)=n^2-n+2.
\]
Third, I cross-checked the VC-dimension conclusions against the critical growth
values:
\[
\growth_{\classA}(4)=16,\ \growth_{\classA}(5)=31,
\qquad
\growth_{\classB}(3)=8,\ \growth_{\classB}(4)=14.
\]
Finally, for Lean, I focused on the mathematically central statements that are
most reusable rather than trying to formalize every combinatorial detail.

The workflow improved over time. Early in the process, AI was most useful for
brainstorming structure and possible proof strategies. Later, I used it more
selectively: mostly to tighten exposition, catch overstatements, and package the
final work into a clean submission format. That made the final pass much more
efficient.

\subsubsection{Conclusion}

AI was useful for organization, drafting, and sanity checks, but the final
correctness of the assignment depended on explicit proofs, independent
verification, and careful revision of any suggestions that were too loose.

\subsubsection{Lean Formalization}

Lean: NO.

\subsubsection{Computational Check}

The computational support mentioned here appears in
\texttt{hw2/code/enumeration\_check.py}, with outputs in \texttt{hw2/figures/}.

\end{solutionbox}

\section{Repository Artifacts and Mapping}

\subsection{Supporting Repository Artifacts}

\begin{itemize}[leftmargin=2.5em]
    \item \texttt{hw2/writeup.tex}: source file for the write-up.
    \item \texttt{hw2/writeup.pdf}: compiled PDF draft.
    \item \texttt{hw2/lean/vc\_dimension\_bound.lean}: Lean target for Problem B1.
    \item \texttt{hw2/lean/quadratic\_embedding.lean}: Lean target for Problem B2.
    \item \texttt{hw2/code/enumeration\_check.py}: small-\(n\) verification for
    Problems A2 and B4.
    \item \texttt{hw2/figures/small\_n\_growth.csv}: table of computed values.
    \item \texttt{hw2/figures/small\_n\_growth.png}: small-\(n\) growth plot.
\end{itemize}

\subsection{Write-Up to Repository Mapping}

\begin{itemize}[leftmargin=2.5em]
    \item Problem A2 \(\rightarrow\)
    \texttt{hw2/code/enumeration\_check.py},
    \texttt{hw2/figures/small\_n\_growth.csv}
    \item Problem B1 \(\rightarrow\)
    \texttt{hw2/lean/vc\_dimension\_bound.lean}
    \item Problem B2 \(\rightarrow\)
    \texttt{hw2/lean/quadratic\_embedding.lean}
    \item Problem B4 \(\rightarrow\)
    \texttt{hw2/code/enumeration\_check.py},
    \texttt{hw2/figures/small\_n\_growth.png}
    \item Full submission draft \(\rightarrow\)
    \texttt{hw2/writeup.tex} and \texttt{hw2/writeup.pdf}
\end{itemize}

\subsection{Lean Targets}

The most appropriate Lean targets in this assignment are:
\begin{itemize}[leftmargin=2.5em]
    \item the general linear-representation VC-dimension upper bound from
    Problem B1,
    \item the quadratic embedding identity from Problem B2.
\end{itemize}
These are the cleanest and most reusable formal statements in the assignment.

\end{document}

\subsection{Problem A2: Growth Function}

\textbf{Question.}
\newline
Using the characterization from Problem A1, compute the growth function
\[
\growth_{\classA}(n),
\]
the number of distinct labelings of \(n\) ordered points realizable by
\(\classA\).

You must count separately:
\begin{enumerate}[leftmargin=2.5em]
    \item no runs of \(1\)s,
    \item exactly one run of \(1\)s,
    \item exactly two runs of \(1\)s.
\end{enumerate}

\begin{solutionbox}

\subsubsection{Restatement}

We want an exact formula for the number of binary strings of length \(n\) with
at most two runs of \(1\)s.

\subsubsection{Explanation}

Problem A1 turns the growth-function problem into a counting problem. Since the
three cases above are disjoint and exhaustive, we can count them separately and
add the results without overcounting.

\subsubsection{Solution Plan}

\begin{enumerate}[leftmargin=2.5em]
    \item Count the all-zero labeling.
    \item Count one-run labelings by choosing the first and last index of the
    run.
    \item Count two-run labelings by choosing four boundary positions.
    \item Add the three counts and simplify.
\end{enumerate}

\subsubsection{Formal Solution / Proof}

\begin{claim}
For every \(n\ge 1\),
\[
\growth_{\classA}(n)=1+\binom{n+1}{2}+\binom{n+1}{4}.
\]
\end{claim}

\begin{proof}
\leavevmode\\
By Problem A1, realizable strings are exactly those with at most two runs of
\(1\)s.

\textbf{Case 1: no runs of \(1\)s.}

There is exactly one such string:
\[
00\cdots0.
\]
So this contributes \(1\).

\textbf{Case 2: exactly one run of \(1\)s.}

Such a string is determined by choosing indices \(a\le b\) and declaring
\[
y_i=
\begin{cases}
1, & a\le i\le b,\\
0, & \text{otherwise}.
\end{cases}
\]
Hence the number of one-run strings is the number of pairs
\[
1\le a\le b\le n,
\]
which is
\[
\binom{n+1}{2}.
\]

\textbf{Case 3: exactly two runs of \(1\)s.}

A two-run string has the form
\[
0\cdots 0\,1\cdots 1\,0\cdots 0\,1\cdots 1\,0\cdots 0.
\]
It is uniquely specified by choosing four boundary positions
\[
0\le s_1<t_1<s_2<t_2\le n
\]
from the \(n+1\) gaps before, between, and after the sample points. The first
run starts after boundary \(s_1\) and ends at boundary \(t_1\); the second run
starts after boundary \(s_2\) and ends at boundary \(t_2\). This gives a
bijection between two-run strings and \(4\)-element subsets of the \(n+1\)
boundary positions. Therefore the number of two-run strings is
\[
\binom{n+1}{4}.
\]

Adding the three disjoint cases yields
\[
\growth_{\classA}(n)=1+\binom{n+1}{2}+\binom{n+1}{4}.
\]
\end{proof}

\subsubsection{Conclusion}

The exact growth function is
\[
\growth_{\classA}(n)=1+\binom{n+1}{2}+\binom{n+1}{4}.
\]

\subsubsection{Lean Formalization}

Lean: OPTIONAL.

\subsubsection{Computational Check}

For \(n=1,\dots,6\), the script verifies the values
\[
2,\ 4,\ 8,\ 16,\ 31,\ 57.
\]

\end{solutionbox}
